\clearpage
\section{Dynamic Ablation and Strike Diagnostics}\label{sec:supp_ablation}

The primary ten-transition endpoint compared liquidation outcomes across IV variants on shared stock paths (Table~\ref{tab:dynamic_ablation}). Mean absolute mark differences measured how much the prices changed along individual paths. Paired mean P\&L differences showed how those changes combined across the ensemble. The three evaluation seeds provided separate realizations of the comparison under the same fitted models and pilot constants (Table~\ref{tab:ablation_seeds}). Expected shortfalls summarized the simulated lower tails of P\&L. The paired standard errors applied to mean differences and did not supply uncertainty intervals for expected shortfall. All reported uncertainty was conditional on the fitted surface and pilot normalization constants. Variation from refitting either component remained outside these estimates.

\begin{table}[H]
\centering
\caption{Liquidation after ten trading transitions (May 12, 2026), with 17 calendar DTE. Each row uses 3{,}000 paths. All monetary quantities are dollars per share. $\Delta$P\&L and mean $|\Delta P|$ compare the variant with the direct-surface reference on identical paths. Parentheses contain the paired Monte Carlo standard error of mean $\Delta$P\&L. More negative 5\% expected shortfall indicates a worse simulated lower tail.}\label{tab:dynamic_ablation}
\small
\begin{tabular}{llrrrr}
\toprule
Contract & IV variant & Mean P\&L & ES$_{5\%}$ & $\Delta$P\&L (SE) & Mean $|\Delta P|$ \\
\midrule
GS put & Frozen IV & 0.67 & -63.08 & +0.28 (0.01) & 0.51 \\
GS put & Direct surface & 0.39 & -62.59 & +0.00 (0.00) & 0.00 \\
GS put & Mean reversion & 0.64 & -62.95 & +0.25 (0.01) & 0.42 \\
GS put & Noise, $\rho=0$ & 0.70 & -62.98 & +0.31 (0.04) & 1.75 \\
GS put & Full factor & 0.34 & -64.60 & -0.05 (0.04) & 1.77 \\
\midrule
GS call & Frozen IV & -0.03 & -71.07 & +0.87 (0.01) & 0.87 \\
GS call & Direct surface & -0.89 & -72.76 & +0.00 (0.00) & 0.00 \\
GS call & Mean reversion & -0.21 & -71.42 & +0.69 (0.01) & 0.69 \\
GS call & Noise, $\rho=0$ & -0.12 & -71.53 & +0.77 (0.04) & 1.93 \\
GS call & Full factor & 0.19 & -69.82 & +1.09 (0.04) & 1.96 \\
\midrule
LLY put & Frozen IV & 7.15 & -35.54 & +1.33 (0.01) & 1.33 \\
LLY put & Direct surface & 5.83 & -36.17 & +0.00 (0.00) & 0.00 \\
LLY put & Mean reversion & 6.96 & -35.53 & +1.13 (0.01) & 1.13 \\
LLY put & Noise, $\rho=0$ & 7.01 & -35.70 & +1.18 (0.04) & 2.04 \\
LLY put & Full factor & 6.66 & -38.27 & +0.83 (0.04) & 1.96 \\
\midrule
LLY call & Frozen IV & 6.31 & -47.30 & +2.23 (0.01) & 2.23 \\
LLY call & Direct surface & 4.08 & -50.42 & +0.00 (0.00) & 0.00 \\
LLY call & Mean reversion & 5.90 & -47.95 & +1.82 (0.00) & 1.82 \\
LLY call & Noise, $\rho=0$ & 5.95 & -48.48 & +1.87 (0.04) & 2.36 \\
LLY call & Full factor & 6.26 & -45.18 & +2.18 (0.04) & 2.41 \\
\midrule
\end{tabular}

\end{table}
\FloatBarrier

\begin{table}[H]
\centering
\caption{Full factor versus direct surface after ten trading transitions, by evaluation seed. Each row uses 1{,}000 paths. Quantities are dollars per share; the paired standard error applies to the mean P\&L difference.}\label{tab:ablation_seeds}
\begin{tabular}{llrrr}
\toprule
Contract & Seed & Mean $|\Delta P|$ & Mean $\Delta$P\&L & Paired SE \\
\midrule
GS put & 20260429 & 1.81 & -0.08 & 0.07 \\
GS put & 20260430 & 1.76 & -0.00 & 0.07 \\
GS put & 20260431 & 1.73 & -0.06 & 0.07 \\
GS call & 20260429 & 1.97 & +1.07 & 0.08 \\
GS call & 20260430 & 1.99 & +1.17 & 0.08 \\
GS call & 20260431 & 1.92 & +1.03 & 0.08 \\
LLY put & 20260429 & 2.01 & +0.83 & 0.07 \\
LLY put & 20260430 & 1.95 & +0.85 & 0.07 \\
LLY put & 20260431 & 1.90 & +0.81 & 0.07 \\
LLY call & 20260429 & 2.44 & +2.21 & 0.07 \\
LLY call & 20260430 & 2.43 & +2.22 & 0.07 \\
LLY call & 20260431 & 2.37 & +2.12 & 0.07 \\
\bottomrule
\end{tabular}

\end{table}

Table~\ref{tab:strike_audit} reports the necessary strike-condition checks at 401 steps. The grid had eleven fixed strikes at the scenario expiry, with fifty stock paths evaluated after ten and twenty transitions plus one shared entry state. As spot moved, some fixed strikes left the calibration moneyness interval; the audit therefore included extrapolated states. Each variant and ticker supplied 2{,}222 bound checks, 2{,}020 adjacent-strike checks, and 1{,}818 convexity checks across calls and puts. The companion machine-readable output reports both lattice depths, each test location, and the violation magnitude. Passing this finite grid does not establish the absence of arbitrage elsewhere or consistency across maturities.

\begin{table}[H]
\centering
\caption{Strike-condition violations exceeding \$0.01 per share at 401 lattice steps. Entries show violations/tests. Vertical denotes the upper bound on an adjacent-strike option-price difference. Max reports the largest positive deviation across the four checks, including deviations below the reporting tolerance.}\label{tab:strike_audit}
\resizebox{\textwidth}{!}{\begin{tabular}{llrrrrr}
\toprule
Ticker & Variant & Bounds & Monotonicity & Vertical & Convexity & Max (\$) \\
\midrule
GS & Frozen IV & 0/2222 & 0/2020 & 0/2020 & 0/1818 & 0.008 \\
GS & Direct surface & 0/2222 & 68/2020 & 67/2020 & 178/1818 & 0.113 \\
GS & Mean reversion & 0/2222 & 2/2020 & 0/2020 & 0/1818 & 0.014 \\
GS & Noise, $\rho=0$ & 0/2222 & 3/2020 & 0/2020 & 0/1818 & 0.014 \\
GS & Full factor & 0/2222 & 3/2020 & 0/2020 & 0/1818 & 0.013 \\
LLY & Frozen IV & 0/2222 & 0/2020 & 0/2020 & 0/1818 & 0.000 \\
LLY & Direct surface & 0/2222 & 0/2020 & 0/2020 & 0/1818 & 0.000 \\
LLY & Mean reversion & 0/2222 & 0/2020 & 0/2020 & 0/1818 & 0.000 \\
LLY & Noise, $\rho=0$ & 0/2222 & 0/2020 & 0/2020 & 0/1818 & 0.000 \\
LLY & Full factor & 0/2222 & 0/2020 & 0/2020 & 0/1818 & 0.000 \\
\bottomrule
\end{tabular}
}
\end{table}
\FloatBarrier
